\chapter{Printing and Streams}\label{ch:printing}

\index{printing}
The 1992 manual says remarkably little about output. Its examples print
values, its chapter on conditions passes streams to functions and calls
\code{format}, but it defines no printing functions and no streams at
all. Both implementations fill the gap, and they fill it differently.
They share a small common core, \code{print}, \code{display},
\code{write-line}, \code{newline} and \code{pp}, which Thomas provided in
1992 and which are enough for most programs. The compiler to C adds
streams, \code{format}, \code{write}, and a generic function,
\code{print-object}, through which a program decides how its own objects
print.

This chapter describes the common core first, then how each kind of
object prints, then the compiler's streams and \code{format}, and
finally what the interactive implementation offers instead.

\section{Two ways to print an object}\label{sec:printing-escape}

\index{escape (printing)}
Every printing function prints an object in one of two styles.

\begin{itemize}
\item \term{Escaped} output, the style of \code{print} and \code{write},
  is meant to be read by a programmer, and where possible by the Dylan
  reader. Strings appear in double quotes with \code{"} and \code{\textbackslash}
  escaped by a backslash, and characters appear as \code{\#\textbackslash a}
  or \code{\#\textbackslash space}.
\item \term{Plain} output, the style of \code{display}, is meant for the
  person running the program. Strings and characters appear as
  themselves, with no quotes or \code{\#\textbackslash}.
\end{itemize}

\noindent
The two styles differ only for strings and characters, and for
collections and objects that contain them. A number, a symbol or a list
of numbers looks the same either way.

\section{The common printing functions}\label{sec:printing-functions}

These five functions exist in both implementations. Each answers its
argument (\code{newline} answers \code{\#f}), so a \code{print} can be
wrapped around an expression without changing its value.

\begin{description}
\item[\code{(print \meta{object})}]\index{print} prints the object
  escaped, on a line of its own.
\item[\code{(write-line \meta{object})}]\index{write-line} does the
  same as \code{print}.
\item[\code{(display \meta{object})}]\index{display} prints the object
  plain, and nothing else: no newline, no space.
\item[\code{(newline)}]\index{newline} ends the line.
\item[\code{(pp \meta{object})}]\index{pp} prints the object escaped on
  a line of its own; in the interactive implementation it is
  \term{pretty-printed}, broken over lines and indented as Dylan code is.
\end{description}

\begin{dylan}
(print "a \"quoted\" string")     ; prints "a \"quoted\" string"
(display "a \"quoted\" string")   ; prints a "quoted" string
(print #\a)                       ; prints #\a
(display #\a)                     ; prints a
(display 42) (newline)            ; prints 42 and ends the line
\end{dylan}

\begin{warnbox}
The two implementations put the newline of \code{print} and
\code{write-line} in different places. The compiler to C prints the
object and \emph{then} a newline, as most languages do. The interactive
implementation, following Thomas, prints a newline \emph{first} and then
the object, leaving the line unfinished. A sequence of \code{print}s
therefore gives the same lines in both, but output that mixes
\code{print} with \code{display} does not:
\begin{dylan}
(print 1) (display "x") (newline)
;; compiler to C:          1
;;                         x
;; interactive:            (an empty line)
;;                         1x
\end{dylan}
(At the read-eval-print loop the \code{;Value:} lines come in between;
this is the output of a compiled file.)
A program meant to run in both should build its lines with
\code{display} and \code{newline}, which behave identically.
\end{warnbox}

In the interactive implementation each of the five takes exactly one
argument (\code{newline} none). In the compiler to C each also takes a
\code{stream:} keyword argument (Section~\ref{sec:printing-streams}), and
\code{print} takes a \code{verbose:} keyword, which it accepts and
ignores. The compiler's \code{pp} does no pretty-printing; it is
\code{print} under another name.

\section{How objects print}\label{sec:printing-classes}

The table below shows how each kind of object prints in each
implementation; ``the same'' means that the interactive implementation
prints exactly what the compiler to C does. Where they differ, the
difference is only in the printed form; the objects are the same.

{\frenchspacing
\begin{longtable}{@{}>{\raggedright\arraybackslash}p{0.27\textwidth}>{\raggedright\arraybackslash}p{0.33\textwidth}>{\raggedright\arraybackslash}p{0.33\textwidth}@{}}
\toprule
Object & Compiler to C & Interactive \\
\midrule
\endhead
\code{42}, \code{-7} & \code{42}, \code{-7} & the same \\
\code{3/4} & \code{3/4} & the same \\
\code{1.5}, \code{9.0}, \code{0.1} & \code{1.5}, \code{9.0}, \code{0.1} & the same \\
\code{1.0e20}, \code{(sqrt 2)} & \code{1.0e20}, \code{1.4142135623730951} & the same \\
\code{\#t}, \code{\#f}, \code{'()} & \code{\#t}, \code{\#f}, \code{()} & the same \\
\code{'foo}, \code{'Hello} & \code{foo}, \code{Hello} & the same \\
\code{foo:} & \code{foo:} & the same \\
\code{'(1 (2 3) . 4)} & \code{(1 (2 3) . 4)} & the same \\
\code{'{}'x} & \code{(quote x)} & \code{'x} \\
\code{(vector 1 "two" \#\textbackslash 3)} & \code{\#(1 "two" \#\textbackslash 3)} & the same \\
a stretchy vector of 1 and 2 & \code{\{<stretchy-vector> 1 2\}} & \code{\#[instance-of <stretchy-vector>]} \\
a deque of 1 and 2 & \code{\{<deque> 1 2\}} & \code{\#[instance-of <deque>]} \\
\code{(range from: 0 up-to: 5)} & \code{\{range 0 by 1 size 5\}} & \code{\#[instance-of <range>]} \\
\code{(range from: 0)} & \code{\{range 0 by 1\}} & \code{\#[instance-of <range>]} \\
a table mapping \code{a} to 1 & \code{\{table a => 1\}} & \code{\#[instance-of <table>]} \\
an array & \code{\{instance of <\%general-array>\}} & \code{\#[instance-of <array>]} \\
the class \code{<point>} & \code{\{the class <point>\}} & \code{\#[dylan-class <point>]} \\
an instance of \code{<point>} & \code{\{instance of <point>\}} & \code{\#[instance-of <point>]} \\
a generic function & \code{\{the generic function print\}} & \code{\#[generic function area]} \\
a method & \code{\{method\}} & \code{\#[method]} \\
\code{(singleton 3)} & \code{\{singleton 3\}} & \code{\#[dylan-singleton 3]} \\
a \code{<simple-error>} & its message, \code{bad thing} & the same \\
\bottomrule
\end{longtable}}

\noindent
Some details:

\begin{itemize}
\item \textbf{Floats} print in the shortest form that reads back as the
  same number, with at least one digit after the point (\code{9.0},
  never \code{9.}). Very large and very small values take an exponent,
  but the implementations switch at different points. The compiler to C
  writes a float without an exponent when its magnitude is at least
  $10^{-5}$ and below $10^{16}$; the interactive implementation, which
  uses Common Lisp's printer, only when it is at least $10^{-3}$ and below
  $10^{7}$. So \code{123456789.0} prints as \code{123456789.0} in the
  first and \code{1.23456789e8} in the second, and \code{0.0001} as
  \code{0.0001} and \code{1.0e-4}.
\item \textbf{Ratios} print as \meta{numerator}\code{/}\meta{denominator}
  in lowest terms; \code{(/ 6 4)} prints as \code{3/2}.
\item \textbf{Symbols} print as they were first spelled. Dylan symbols
  are not case-sensitive, but the reader remembers the spelling a symbol
  was first entered in, so \code{'Hello} prints as \code{Hello}
  \dylanbook{chapter 15}. \textbf{Keywords} print with their colon.
\item \textbf{Characters} with names print escaped by name:
  \code{\#\textbackslash space}, \code{\#\textbackslash newline},
  \code{\#\textbackslash tab}, \code{\#\textbackslash return},
  \code{\#\textbackslash page}, \code{\#\textbackslash null},
  \code{\#\textbackslash rubout}, \code{\#\textbackslash backspace}.
\item \textbf{Lists and vectors} print their elements in the same style
  as the whole: \code{(display (list "a" \#\textbackslash b))} prints
  \code{(a b)}.
\item \textbf{Conditions} print as their message, the format string
  filled in with the format arguments, as \dylanbook{chapter 13}
  suggests. A \code{<type-error>} prints as
  ``\meta{value} is not an instance of \meta{type}'', where the type is
  printed as a class prints: \code{3 is not an instance of \{the class
  <string>\}} in the compiler to C, \code{3 is not an instance of
  <string>} in the interactive implementation.
\item \textbf{Functions} that the language predefines as plain methods,
  such as \code{car}, print as methods. Which predefined functions are
  generic differs between the implementations: \code{print} is a generic
  function in the compiler to C and a method in the interactive one.
\end{itemize}

\section{Customizing printing with \code{print-object}}\label{sec:printing-print-object}

\index{print-object}
In the compiler to C every printing function, \code{format} included,
works by calling one generic function:

\begin{dylan}
(print-object object stream escape?)
\end{dylan}

\noindent
\code{print-object} writes \meta{object} to \meta{stream}; \meta{escape?}
is \code{\#t} for escaped output (\code{print}, \code{write}, \code{\textasciitilde s})
and \code{\#f} for plain output (\code{display}, \code{\textasciitilde a}). The library
defines a method for each of its own classes, and a default method on
\code{<object>} that prints \code{\{instance of \meta{class}\}}. To make
the instances of your own class print differently, add a method:

\begin{dylan}
(define-class <point> (<object>)
  (x init-keyword: x: init-value: 0)
  (y init-keyword: y: init-value: 0))

(define-method print-object ((p <point>) stream escape?)
  (format stream "{point ~s, ~s}" (x p) (y p)))

(define p (make <point> x: 3 y: 4))
(print p)                     ; prints {point 3, 4}
(print (list p p))            ; prints ({point 3, 4} {point 3, 4})
(format #t "p is ~a~%" p)     ; prints p is {point 3, 4}
\end{dylan}

\noindent
The method is used wherever a point is printed, including inside lists,
vectors and other collections, and by \code{format}. Write the output
to the \meta{stream} argument, never to \code{*standard-output*}
directly, or printing into a string will not work.

A method can honor the \meta{escape?} argument, so that \code{print}
shows something a programmer can recognize and \code{display} shows what
a user should see:

\begin{dylan}
(define-class <name> (<object>) (text init-keyword: text:))

(define-method print-object ((n <name>) stream escape?)
  (if escape?
      (format stream "#<name ~s>" (text n))
      (display (text n) stream: stream)))

(define n (make <name> text: "Ada"))
(print n)                       ; prints #<name "Ada">
(display n)                     ; prints Ada
(format #t "~a / ~s~%" n n)     ; prints Ada / #<name "Ada">
\end{dylan}

\noindent
Methods are inherited and can call \code{next-method} like any other:

\begin{dylan}
(define-class <temp> (<object>) (deg init-keyword: deg:))
(define-method print-object ((t <temp>) stream escape?)
  (display (deg t) stream: stream)
  (display " degrees" stream: stream))

(define-class <hot> (<temp>))
(define-method print-object ((t <hot>) stream escape?)
  (display "hot! " stream: stream)
  (next-method))

(print (make <hot> deg: 99))    ; prints hot! 99 degrees
\end{dylan}

\noindent
\code{print-object} is an extension: the 1992 manual has no such
function, and the interactive implementation does not provide one
(Section~\ref{sec:printing-lisp}).

\section{Streams in the compiler to C}\label{sec:printing-streams}

\index{stream}
The compiler to C provides a small set of streams, enough for the uses
the 1992 manual makes of them. They are an extension, described with the
compiler's other extensions in Chapter~\ref{ch:extensions}.

\begin{description}
\item[\code{<stream>}]\index{<stream>@\code{<stream>}} The class of all
  streams.
\item[\code{*standard-output*}]\index{*standard-output*@\code{*standard-output*}}
  A stream that writes to the program's standard output.
\item[\code{*standard-error*}]\index{*standard-error*@\code{*standard-error*}}
  A stream that writes to the program's standard error.
\item[\code{<string-stream>}]\index{<string-stream>@\code{<string-stream>}}
  A stream that collects what is written to it, in memory. You make
  one by calling \code{make}: \code{(make <string-stream>)}.
\item[\code{(stream-contents \meta{string-stream})}]\index{stream-contents}
  Answers a new string of everything written to the stream so far. The
  stream keeps its contents; calling \code{stream-contents} again answers
  them again, with anything written since.
\end{description}

\noindent
Every printing function takes a \code{stream:} keyword argument naming
the stream to write to. It defaults to \code{*standard-output*}, and
\code{\#t} (or \code{\#f}) also means \code{*standard-output*}.

\begin{dylan}
(define s (make <string-stream>))
(print 42 stream: s)
(display "hello" stream: s)
(write #\x stream: s)
(display (stream-contents s))
;; prints  42
;;         hello#\x

(display "warning: disk nearly full" stream: *standard-error*)
(newline stream: *standard-error*)
\end{dylan}

\noindent
\code{write}\index{write} prints its argument escaped, like
\code{print} but with no newline:

\begin{description}
\item[\code{(write \meta{object} \#key stream)}] prints the object
  escaped and answers it.
\end{description}

\noindent
\code{<stream>} cannot usefully be subclassed by a program. The
operation that puts characters on a stream is internal to the library,
so printing to an instance of your own subclass of \code{<stream>} fails
with ``no applicable method''.

\section{\code{format}}\label{sec:printing-format}

\index{format}
\begin{dylan}
(format stream control-string argument ...)
\end{dylan}

\noindent
\code{format} writes \meta{control-string} to \meta{stream}, replacing
each \term{directive}, a tilde and a letter, by output for the next
argument. Its directives are a small subset of Common Lisp's, the ones
the 1992 manual's examples use:

\begin{center}
\begin{tabular}{@{}ll@{}}
\toprule
Directive & Output \\
\midrule
\code{\textasciitilde a} & the next argument, plain, as \code{display} prints it \\
\code{\textasciitilde s} & the next argument, escaped, as \code{print} and \code{write} print it \\
\code{\textasciitilde d} & the next argument, plain (the same as \code{\textasciitilde a}) \\
\code{\textasciitilde\%} & a newline \\
\code{\textasciitilde\&} & a newline (the same as \code{\textasciitilde\%}) \\
\code{\textasciitilde\textasciitilde} & a tilde \\
\bottomrule
\end{tabular}
\end{center}

\noindent
The letters may be upper or lower case. Any other tilde sequence is
copied to the output unchanged. There are no field widths, no
padding and no number bases: \code{\textasciitilde d} prints an integer
exactly as \code{display} would, and \code{\textasciitilde\&} always
starts a new line, even at the beginning of one.

The \meta{stream} argument is a stream, or one of two booleans:

\begin{itemize}
\item \code{\#t} writes to \code{*standard-output*}, and \code{format}
  answers \code{\#f}.
\item \code{\#f} writes to a new string, and \code{format} answers the
  string.
\end{itemize}

\begin{dylan}
(format #t "~a + ~a = ~a~%" 1 2 (+ 1 2))    ; prints 1 + 2 = 3
(format #f "~a + ~a = ~a" 1 2 (+ 1 2))      ; => "1 + 2 = 3"
(format #f "~A ~S ~D ~z" "x" "x" 5)         ; => "x \"x\" 5 ~z"
(format *standard-error* "to stderr~%")
\end{dylan}

\noindent
Too few arguments for the directives is an error
(``\code{format: too few arguments for "\textasciitilde a and \textasciitilde a"}''); arguments left
over after the last directive are ignored.

\section{Printing in the interactive implementation}\label{sec:printing-lisp}

\index{interactive implementation!printing}
The interactive implementation has exactly the five functions of
Section~\ref{sec:printing-functions}, each with one argument and no
\code{stream:} keyword. It has no streams, no \code{format}, no
\code{write} and no \code{print-object}, so the way an object prints
cannot be changed from Dylan: an instance of a class of your own always
prints as \code{\#[instance-of \meta{class}]}. To show its contents,
print them yourself:

\begin{dylan}
(define-method show ((p <point>))
  (display "{point ") (display (x p))
  (display ", ") (display (y p)) (display "}"))
\end{dylan}

\noindent
The read-eval-print loop prints each result escaped after
\code{;Value:}, so at the \code{?} prompt you seldom need \code{print}
at all (Chapter~\ref{ch:lisp}).

\subsection{Pretty-printing}

\index{pp}
\code{pp} uses the Common Lisp pretty printer, set up to show Dylan's
notation, to lay out a long list as code:

\begin{repl}
? (pp '(define-method foo ((x <integer>))
         (bind ((y (* x x)) (z (+ x x x x x x x x x x x x x x x x x x x x)))
           (list y z y z y z))))

(define-method foo ((x <integer>))
 (bind ((y (* x x)) (z (+ x x x x x x x x x x x x x x x x x x x x)))
  (list y z y z y z)))
\end{repl}

\subsection{Formatted output through Common Lisp}

\index{lisp-procedure}
When you need formatted output in the interactive implementation, Common
Lisp's own \code{format} can be reached with \code{lisp-procedure}
(Section~\ref{sec:lisp-interop}). It has every Common Lisp directive,
field widths included. Pass \code{\#t} for standard output, or the empty
list, which is Lisp's \code{NIL}, to get a string back:

\begin{repl}
? (define cl-format (lisp-procedure 'format))
;Value: cl-format

? (cl-format #t "~5,2f|~8d|~r~%" 3.14159 42 12)
 3.14|      42|twelve
;Value: ()

? (cl-format '() "~a-~a" 1 2)
;Value: "1-2"
\end{repl}

\begin{warnbox}
Common Lisp's \code{format} prints Dylan values as Common Lisp sees
them, not as Dylan's printer does: \code{\#t} comes out as \code{T}, a
keyword \code{foo:} as \code{FOO}, and \code{\#f}, which is not a Lisp
value at all, as \code{\#f} only because the implementation teaches Lisp
to print it so. Use it for numbers and strings. Do not pass \code{\#f}
as the destination: to Common Lisp it is an ordinary true object, not
\code{NIL}.
\end{warnbox}

\section{Writing portable output}\label{sec:printing-portable}

A program that must run in both implementations should

\begin{itemize}
\item build its output from \code{display} and \code{newline}, which
  behave the same in both;
\item use \code{print} only for whole lines, remembering where each
  implementation puts the newline;
\item avoid \code{format}, \code{write}, streams and \code{print-object},
  which only the compiler to C has; and
\item not depend on the printed form of collections other than lists,
  vectors and strings, or of classes, functions and instances, which
  differ (Section~\ref{sec:printing-classes}).
\end{itemize}
